% =====================================================================
% FLAIR --- NEJM AI Datasets, Benchmarks, and Protocols
% Manuscript outline (working document, not the submission template)
%
% Target venue:  NEJM AI, "Datasets, Benchmarks, and Protocols" (AIdbp)
% Envelope:      <= 3000 main-text words; <= 5 figures+tables combined;
%                <= 40 references; <= 10 authors; structured abstract.
% Exemplars (in docs/paper/lit/):
%   - AIdbp2500418  Verifying Facts in Patient Care Documents (VeriFact)
%   - AIdbp2500516  EchoNext-Mini Dataset and Baseline AI Model
%   - AIdbp2401267  Pragmatic Trial Operations Playbook (protocol-flavored)
% =====================================================================

\documentclass[11pt]{article}

\usepackage[margin=1in]{geometry}
\usepackage[T1]{fontenc}
\usepackage{lmodern}
\usepackage{microtype}
\usepackage{amsmath, amssymb}
\usepackage{booktabs}
\usepackage{tabularx}
\usepackage{array}
\usepackage{enumitem}
\usepackage[hidelinks]{hyperref}
\usepackage{xcolor}
\usepackage[round]{natbib}

\setlength{\parskip}{0.4em}
\setlength{\parindent}{0pt}

% Editorial-note macro: word/figure budgets, drafting prompts, TODOs
\definecolor{notecolor}{RGB}{120,120,120}
\newcommand{\note}[1]{{\color{notecolor}\textit{[#1]}}}
\newcommand{\todo}[1]{{\color{red}\textbf{TODO:} #1}}

\title{\textbf{FLAIR: A Federated Benchmark and Baseline Repository for
ICU Prediction Across 12+ Academic Medical Centers Using the Common
Longitudinal ICU Format} \\[0.4em]
\large Manuscript outline --- NEJM AI Datasets, Benchmarks, and Protocols}
\author{CLIF Consortium \\ FLAIR Working Group}
\date{\today}

\begin{document}
\maketitle

\note{Working outline. Bulleted content cues sit under every section so
drafting prose is paint-by-numbers. Hard envelope: 3000 main-text words,
5 figures+tables combined, 40 references, structured abstract.}

% ---------------------------------------------------------------------
\section*{Front matter}
% ---------------------------------------------------------------------
\begin{itemize}[leftmargin=1.5em]
  \item \textbf{Title (working):} FLAIR: A Federated Benchmark for ICU
        Prediction Across 12+ Academic Medical Centers Using the Common
        Longitudinal ICU Format.
  \item \textbf{Short title} ($\le$ 50 char): FLAIR: Federated ICU Benchmark.
  \item \textbf{Authors / affiliations:} \todo{populate; up to 10 authors;
        flag equal-contribution and co-senior authorship in a title-page
        footnote. ORCIDs required.}
  \item \textbf{Corresponding author:} \todo{name, email, institution.}
  \item \textbf{Funding:} \todo{grant numbers; restated as a "Funded by"
        line at the end of the abstract.}
  \item \textbf{Conflicts of interest:} \todo{ICMJE + Convey forms within
        12 months for every author.}
  \item \textbf{AI-use disclosure:} \todo{name LLM(s); state which sections
        or artifacts they helped produce. Restate in cover letter.}
  \item \textbf{Word / figure / table / reference counts:}
        Main text \note{$\le$ 3000}; figures+tables \note{$\le$ 5};
        references \note{$\le$ 40}.
\end{itemize}

% ---------------------------------------------------------------------
\section*{Structured abstract \note{$\approx$ 250 words}}
% ---------------------------------------------------------------------

\subsubsection*{Background}
\begin{itemize}[leftmargin=1.5em]
  \item ICU prediction models trained on public datasets (MIMIC, eICU,
        HiRID) routinely fail to generalize beyond the institutions that
        produced their training data.
  \item Patient-level data cannot be centralized under HIPAA / IRB
        constraints, so multi-institutional external validation requires
        a federated evaluation framework.
  \item No prior ICU benchmark combines a shared data schema, auditable
        per-task code, aggregate-only outputs, and TRIPOD-AI reporting.
\end{itemize}

\subsubsection*{Methods}
\begin{itemize}[leftmargin=1.5em]
  \item Federated benchmark built atop CLIF v2.1 across N CLIF-consortium
        sites \note{fill N at submission}.
  \item Five clinically grounded tasks: T1 ICU daily mortality, T2 ICU
        daily LTACH discharge, T3 extubation failure $<$24h, T4 sepsis
        (antibiotic-anchored) $<$6h, T5 AKI KDIGO-2 $<$24h.
  \item Reference baseline kit: (i) a generic CLIF event tokenizer that
        maps every clinical event into a token and emits per-window
        count features, and (ii) two reference models --- L2-penalized
        logistic regression and XGBoost --- trained on MIMIC-CLIF and
        shipped inside the CLIF-WorkBench container for federated
        evaluation. Sites emit four aggregated JSON reports
        (\texttt{table1}, \texttt{tripod\_ai}, \texttt{fairness},
        \texttt{curves}).
\end{itemize}

\subsubsection*{Results}
\begin{itemize}[leftmargin=1.5em]
  \item Combined cohort: \todo{n hospitalizations, n ICU stays, n sites
        across the consortium}. MIMIC-CLIF baseline cohort (development
        entry point): 84{,}854 ICU stays at the broadest task (T4 sepsis);
        12{,}438--65{,}577 stays for the four task-specific cohorts
        (T1/T2 daily disposition, T3 extubation, T5 AKI); per-task
        positive-rate range 1.7\%--57.4\% at decision-row granularity
        (see Table~\ref{tab:mimic-snapshot}).
  \item Per-task AUROC / AUPRC (median, IQR across sites) \todo{fill from
        federated runs}; calibration drift quantified between MIMIC-CLIF
        and consortium sites.
  \item Fairness deltas across sex, race, ethnicity, and age bins; bit-
        identical SHA-256 hashes across sites confirm protocol fidelity.
\end{itemize}

\subsubsection*{Conclusions}
\begin{itemize}[leftmargin=1.5em]
  \item FLAIR enables external validation of ICU AI models without data
        sharing. Code, container, and per-task schematics released openly.
        (Funded by \todo{...}.)
\end{itemize}

% ---------------------------------------------------------------------
\section{Introduction \note{$\approx$ 600 words}}
% ---------------------------------------------------------------------
\begin{itemize}[leftmargin=1.5em]
  \item \textbf{P1 --- Clinical stakes.} ICU decisions are real-time and
        high-consequence; rich EHR signal is available but underused.
  \item \textbf{P2 --- The saturation problem in public benchmarks.}
        YAIB~\citep{vandewater2023yaib} and HiRID-ICU-Benchmark
        \citep{yeche2021hirid} harmonize public datasets (MIMIC, eICU,
        HiRID, AUMCdb); strong internal performance does not imply
        external generalization.
  \item \textbf{P3 --- The privacy barrier.} HIPAA / IRB rules block
        centralization of PHI; multi-site validation must therefore be
        federated.
  \item \textbf{P4 --- Why prior federated ICU efforts fall short.} Most
        rely on ad-hoc schemas, lack auditable per-task code, and do not
        emit TRIPOD-AI-aligned reports.
  \item \textbf{P5 --- FLAIR's contribution.} Five clinically grounded
        tasks; CLIF v2.1 schema-locked code; single-file polars task
        scripts stamped with SHA-256; \textbf{a generic event-tokenizer
        ETL plus two reference baselines (logistic regression, XGBoost)
        released as the FLAIR baseline kit}; aggregate-only report
        bundle; baseline + federated results across the CLIF consortium.
  \item \textbf{Figure 1 callout} \note{counts toward 5-item budget}:
        \emph{Summary of Contributions} panel modeled on EchoNext-Mini
        Fig.\ 1 \citep{hughes2026echonextmini} --- one side shows the
        federated workflow (MIMIC-CLIF $\to$ event tokenizer $\to$
        count features $\to$ \{logistic regression, XGBoost\} $\to$
        frozen baseline kit (Singularity) $\to$ consortium sites $\to$
        aggregate JSONs); the other shows the five-task taxonomy.
\end{itemize}

% ---------------------------------------------------------------------
\section{Methods \note{$\approx$ 1200 words --- heaviest section}}
% ---------------------------------------------------------------------

\subsection*{Dataset curation}
\begin{itemize}[leftmargin=1.5em]
  \item CLIF v2.1 data standard maintained by 12+ US academic medical
        centers; combined coverage $\sim$2{,}000 ICU beds nationwide.
  \item MIMIC-CLIF as public development entry point ($\sim$70k ICU stays;
        PhysioNet); identical schema guarantees code runs unmodified on
        consortium data.
  \item Site enrollment and IRB framework: each site IRB-approved
        independently; central protocol coordination at the University
        of Chicago \todo{confirm}.
  \item CLIF tables exposed by \texttt{flair\_benchmark/\_clif.py}:
        \texttt{hospitalization}, \texttt{adt}, \texttt{patient},
        \texttt{vitals}, \texttt{labs}, \texttt{respiratory\_support},
        \texttt{medication\_admin\_intermittent},
        \texttt{medication\_admin\_continuous}.
  \item Data-quality gates inherited from \texttt{clifpy} validators;
        timezone handling per site \texttt{clif\_config.json}.
\end{itemize}

\subsection*{Task definitions}

\paragraph{Taxonomy.}
Each task is declared on three axes in its \texttt{META} dict:
\emph{task class} (outcome / deterioration / intervention),
\emph{temporal pattern} (episode / continuous), and
\emph{prediction unit} (one-per-stay / one-per-timestep).
Full algorithmic specs sit in \texttt{docs/task\_schematics/} and are
included by reference in the Supplement.

\begin{itemize}[leftmargin=1.5em]
  \item \textbf{T1 --- ICU Daily In-Hospital Mortality (7\,AM).}
        one-per-timestep, 24\,h lookback, daily 07:00 local anchor.
        Cohort: first ICU stay; discharge $\in$ \{LTACH, mortality\}.
        \emph{MIMIC-CLIF:} 12{,}438 stays / 65{,}989 daily rows;
        per-row positive rate 57.4\%; stays-ever-positive 67.1\%.
  \item \textbf{T2 --- ICU Daily LTACH Discharge (7\,AM).} Complementary
        label to T1; same cohort and grid.
        \emph{MIMIC-CLIF:} same 12{,}438 stays / 65{,}989 rows;
        per-row positive rate 42.6\%; stays-ever-positive 32.9\%
        (mortality-vs-LTACH disposition is asymmetric).
  \item \textbf{T3 --- Extubation Failure within 24\,h.} one-per-stay,
        anchored at the first confirmed extubation event; consensus-window
        IMV episode detector in \texttt{\_ventilation\_episodes.py}.
        Label = death OR reintubation within 24\,h post-extubation.
        \emph{MIMIC-CLIF:} 26{,}771 extubation episodes;
        prevalence 5.4\% (1{,}456 / 26{,}771).
  \item \textbf{T4 --- Sepsis (antibiotic-anchored) within 6\,h.}
        one-per-timestep, hourly grid; $t_{\text{zero}}$ = first qualifying
        systemic IV antibiotic; label-leakage gate drops decision points
        $\ge t_{\text{zero}}$ \citep{kamran2024sepsis}.
        \emph{MIMIC-CLIF:} 84{,}854 stays / 7{,}093{,}633 hourly rows;
        per-row positive rate 3.3\%; stays-ever-positive 54.9\%.
  \item \textbf{T5 --- AKI KDIGO-2 within 24\,h.} one-per-timestep, hourly
        grid; SCr-only criterion; Koyner-2018 exclusions
        \citep{koyner2018aki}.
        \emph{MIMIC-CLIF:} 65{,}577 stays / 4{,}135{,}079 hourly rows;
        per-row positive rate 1.7\%; stays-ever-positive 4.8\%.
\end{itemize}

\paragraph{MIMIC-CLIF cohort snapshot.} The numbers below come from
\texttt{flair extract-metadata --clif-config clif\_config\_mimic.json}
(\texttt{flair\_benchmark} v2.0.0; HIPAA cell-suppression at $n{<}10$;
SOFA computed over the first 24\,h of the first ICU stay via
\texttt{clifpy}; Charlson and Elixhauser computed from
\texttt{HospitalDiagnosis} where available --- \emph{note: 49--54\% of
stays are missing comorbidity scores in the MIMIC-CLIF export, a known
data-quality gap to flag in the Limitations}).

\begin{table}[h]
\centering
\small
\setlength{\tabcolsep}{4pt}
\begin{tabular}{lrrrrr}
\toprule
                       & \textbf{T1 mort.} & \textbf{T2 LTACH} & \textbf{T3 ext-fail} & \textbf{T4 sepsis} & \textbf{T5 AKI} \\
\midrule
ICU stays              & 12{,}438        & 12{,}438         & 26{,}771           & 84{,}854         & 65{,}577        \\
Decision rows          & 65{,}989        & 65{,}989         & 26{,}771           & 7{,}093{,}633   & 4{,}135{,}079   \\
Positive rate (row)    & 57.4\%          & 42.6\%           & 5.4\%              & 3.3\%            & 1.7\%           \\
Stays ever positive    & 67.1\%          & 32.9\%           & ---                & 54.9\%           & 4.8\%           \\
\midrule
Sex: Female / Male     & 5{,}546 / 6{,}892 & 5{,}546 / 6{,}892 & 10{,}033 / 16{,}738 & 37{,}552 / 47{,}302 & 28{,}867 / 36{,}710 \\
Age 65+ share          & 66.6\%          & 66.6\%           & 54.1\%             & 54.6\%           & 55.1\%          \\
SOFA 24h (median, IQR) & 5 (2--8)        & 5 (2--8)         & 5 (4--8)           & 3 (1--6)         & 3 (1--5)        \\
Charlson (median, IQR) & 4 (2--5)        & 4 (2--5)         & 2 (0--4)           & 2 (0--4)         & 2 (0--4)        \\
\bottomrule
\end{tabular}
\caption{MIMIC-CLIF cohort snapshot per task, generated by
\texttt{flair extract-metadata}. \emph{Decision rows} = one-per-stay for
T3 (extubation episode), one-per-timestep otherwise (daily at 07:00 for
T1/T2; hourly for T4/T5). \emph{Age 65+} share is over distinct ICU
stays. \emph{Stays ever positive} is omitted for T3 because the task is
one-per-stay (the row-level rate is the stay-level rate).}
\label{tab:mimic-snapshot}
\end{table}

\note{This table is a drafting hook for Results---Table~1
and Methods---Cohort. It is not part of the 5-item NEJM AI
figure/table budget; the submission Table~1 will mirror this layout
across consortium sites with site-level median/IQR.}

\paragraph{Splits.} Temporal date-based train/test split; MIMIC fallback
70/30 when split dates not specified.

\paragraph{Cohort parquet schema.} Ten columns:
\texttt{hospitalization\_id}, \texttt{window\_start}, \texttt{window\_end},
\texttt{<label>}, \texttt{split}, \texttt{age\_at\_admission},
\texttt{sex\_category}, \texttt{race\_category},
\texttt{ethnicity\_category}, and the task-specific \texttt{label\_*} column.

\subsection*{Federated protocol}
\begin{itemize}[leftmargin=1.5em]
  \item Two-layer privacy model: (i) CLIF-WorkBench Singularity container
        --- no internet, no root, immutable read-only image; (ii)
        aggregate-only outputs from \texttt{flair build-report}.
  \item Library-enforced HIPAA safeguards: cell suppression in
        \texttt{table1.json} for counts $<$10; subgroup drop in
        \texttt{fairness.json} for subgroups with $n<10$; calibration-bin
        suppression in \texttt{curves.json}.
  \item Audit chain: SHA-256 of each task script is stamped into every
        output JSON; the central aggregator verifies byte-identical task
        code across sites.
  \item Submission contract: user appends \texttt{y\_pred}/\texttt{y\_prob}
        to the cohort parquet; \texttt{flair build-report} emits exactly
        four JSON files. No patient-level data ever leaves the site.
        For the FLAIR baseline kit, the container additionally ships
        the tokenizer and frozen weights; sites invoke
        \texttt{flair predict-baseline <task>} instead of writing their
        own feature/model code.
\end{itemize}

\subsection*{Baseline kit}

\paragraph{Generic event tokenizer.} Each row in a CLIF event table
(\texttt{vitals}, \texttt{labs},
\texttt{medication\_admin\_intermittent},
\texttt{medication\_admin\_continuous}, \texttt{respiratory\_support})
is mapped to a deterministic token of the form
\texttt{<table>.<category>} (e.g., \texttt{vitals.heart\_rate},
\texttt{labs.creatinine}, \texttt{meds.norepinephrine}). For each
cohort window, FLAIR emits a sparse token $\to$ integer-count vector
counting events in $[\texttt{window\_start}, \texttt{window\_end}]$.
The tokenizer is config-free and exhaustive --- every category present
in the input flows through with no per-task curation. Because tokens
are derived from CLIF schema-locked category fields, the same code
runs unchanged on MIMIC-CLIF and consortium sites. Module:
\texttt{flair\_benchmark/baseline/tokenizer.py}.

\paragraph{Logistic-regression baseline.} Sparse count features $\to$
L2-penalized logistic regression (\texttt{sklearn}); submitted as
\texttt{y\_prob} for the binary tasks.

\paragraph{XGBoost baseline.} The same sparse count features $\to$
XGBoost classifier with default depth and learning rate; submitted as
\texttt{y\_prob}.

Both models are trained on MIMIC-CLIF (70/30 temporal split via
\texttt{\_split.py}) for each of T1--T5, serialized into the
CLIF-WorkBench Singularity container, and shipped to each consortium
site. Hyperparameters, the token dictionary, and the missingness
convention (absent token $=$ zero count) are deferred to the Supplement
(Model Cards + Datasheet).

\subsection*{Reporting and statistical analysis}
\begin{itemize}[leftmargin=1.5em]
  \item TRIPOD+AI \citep{collins2024tripodai} compliance via
        \texttt{flair\_benchmark/report/tripod\_ai.py} auto-population.
  \item Primary metrics: AUROC, AUPRC; 95\% CIs via 1000-iteration
        bootstrap.
  \item Calibration: slope, intercept, Brier score, reliability curves
        (downsampled raw points in \texttt{curves.json}).
  \item Fairness: per-subgroup AUROC and parity gaps across sex, race,
        ethnicity, and age bins.
  \item Central dashboard renders ROC / PR / calibration curves from the
        aggregated \texttt{curves.json} without re-accessing patient data.
\end{itemize}

% ---------------------------------------------------------------------
\section{Results \note{$\approx$ 700 words}}
% ---------------------------------------------------------------------

\subsection*{Cohort}
\begin{itemize}[leftmargin=1.5em]
  \item \textbf{Table 1.} Combined cohort characteristics across N sites
        (HIPAA-suppressed cells where applicable), stratified by split
        and label.
  \item \textbf{Figure 2 callout:} CONSORT-style flow diagram per task
        (raw $\to$ cohort filters $\to$ train/test). Mirrors EchoNext-Mini
        Fig.\ 2 layout.
  \item Reported counts: total hospitalizations, total ICU stays per task,
        per-task prevalence, demographic distribution.
\end{itemize}

\subsection*{Baseline performance on MIMIC-CLIF}
\begin{itemize}[leftmargin=1.5em]
  \item Per-task internal test AUROC / AUPRC with 95\% bootstrap CIs,
        reported separately for each model family (logistic regression,
        XGBoost).
  \item Calibration slope / intercept / Brier score per task per model
        family.
\end{itemize}

\subsection*{Federated performance across the CLIF consortium}
\begin{itemize}[leftmargin=1.5em]
  \item \textbf{Table 2.} Per-task site-level summary (median, IQR, range
        for AUROC / AUPRC across N sites), with a
        \texttt{model\_family} column so each task carries one row per
        family per site.
  \item \textbf{Figure 3.} Forest plot of per-site AUROC by task with
        two glyphs per task (logistic regression vs XGBoost); vertical
        reference line at the MIMIC-CLIF internal benchmark for that
        family.
  \item Quantify calibration drift between MIMIC-CLIF and held-out
        consortium sites, separately by model family.
\end{itemize}

\subsection*{Fairness}
\begin{itemize}[leftmargin=1.5em]
  \item \textbf{Figure 4.} Heatmap of per-subgroup AUROC (rows = tasks,
        columns = subgroups) with parity-gap shading. Cells with $n<10$
        are suppressed.
\end{itemize}

\subsection*{Reproducibility evidence}
\begin{itemize}[leftmargin=1.5em]
  \item Bit-identical SHA-256 task-script hashes recorded in every output
        JSON across sites; the central aggregator's audit log confirms
        protocol fidelity.
\end{itemize}

% ---------------------------------------------------------------------
\section{Discussion \note{$\approx$ 400 words}}
% ---------------------------------------------------------------------
\begin{itemize}[leftmargin=1.5em]
  \item \textbf{Generalization.} How much do MIMIC-trained baselines drop
        when evaluated on held-out consortium sites? Quantify and
        contextualize.
  \item \textbf{Positioning vs.\ prior benchmarks.} YAIB
        \citep{vandewater2023yaib} and HiRID-ICU-Benchmark
        \citep{yeche2021hirid} harmonize public sources only; FLAIR's
        differentiator is federation across private data with an audit
        chain.
  \item \textbf{Limitations.}
        \texttt{clifpy} data-quality dependency; T5 CKD-stage exclusion
        not yet implemented (CLIF \texttt{Diagnoses} table not exposed in
        FLAIR loaders); T4 antibiotic drug-list curation is an open
        dependency; baseline tokens are pure counts within the
        prediction window with no time-of-occurrence, magnitude, or
        trend features (deliberate floor; richer encodings are left to
        user submissions); no urine-output criterion for AKI.
  \item \textbf{Bias and fairness.} Consortium sites self-select; rural
        and safety-net hospitals are under-represented.
  \item \textbf{Intended use.} External validation tool for development
        and method comparison --- \emph{not} a leaderboard endorsing
        production deployment.
  \item \textbf{Regulatory framing.} Per FDA CDS guidance and IMDRF SaMD
        clinical-evaluation guidance (see \texttt{docs/paper/lit/}), the
        benchmark targets the ``inform clinical management'' use class,
        not autonomous decision-making.
\end{itemize}

% ---------------------------------------------------------------------
\section{Conclusions \note{$\approx$ 120 words}}
% ---------------------------------------------------------------------
\begin{itemize}[leftmargin=1.5em]
  \item One paragraph. FLAIR delivers a federated, schema-locked,
        audit-traceable ICU benchmark \textbf{paired with a reusable
        baseline kit (event tokenizer + logistic regression + XGBoost)
        that defines the floor of expected performance at every site};
        release includes code, container, baseline weights, and
        per-task LaTeX schematics. Community extension via GitHub PRs
        is invited.
\end{itemize}

% ---------------------------------------------------------------------
\section*{Required statements}
% ---------------------------------------------------------------------
\begin{itemize}[leftmargin=1.5em]
  \item \textbf{Data and Code Availability.} GitHub
        \texttt{clif-consortium/FLAIR} (Apache 2.0); MIMIC-CLIF on
        PhysioNet; baseline weights on Zenodo \todo{DOI placeholder};
        Singularity image tag on GHCR \todo{tag placeholder}.
  \item \textbf{IRB.} Each consortium site obtained independent IRB
        approval; central protocol coordination at the University of
        Chicago \todo{confirm}.
  \item \textbf{Funding.} \todo{grants}.
  \item \textbf{Author Contributions.} CRediT taxonomy slots \todo{fill}.
  \item \textbf{Conflicts of Interest.} \todo{disclosures}.
  \item \textbf{AI-Use Disclosure.} \todo{which LLM helped where}.
\end{itemize}

% ---------------------------------------------------------------------
\section*{Supplementary Appendix \note{no NEJM AI hard cap}}
% ---------------------------------------------------------------------
\begin{itemize}[leftmargin=1.5em]
  \item Per-task LaTeX schematics (T1--T5), compiled from
        \texttt{docs/task\_schematics/}; the sepsis-T4 schematic
        \citep{kamran2024sepsis} already exists and is the template.
  \item Datasheet for Datasets (Gebru-style).
  \item Model Cards (Mitchell-style) for the two baseline models
        (logistic regression and XGBoost).
  \item Tokenizer specification --- deterministic
        \texttt{<table>.<category>} mapping; vocabulary size on
        MIMIC-CLIF; missingness convention (absent token $=$ zero
        count); per-token frequency histograms.
  \item Full TRIPOD+AI checklist.
  \item Hyperparameter table, feature dictionary, missingness table.
  \item Per-site Table-1 details (cells with $n \ge 10$).
\end{itemize}

% ---------------------------------------------------------------------
\bibliographystyle{plainnat}
\bibliography{references}

\end{document}
